% © 2026 Ing. David Jaroš — CC BY-NC-ND 4.0
%
% This work is licensed under a Creative Commons Attribution-NonCommercial-NoDerivatives
% 4.0 International License (CC BY-NC-ND 4.0).
%
% License History: Earlier drafts (up to v0.3) were released under CC BY 4.0.
% From v0.4 onward, all material is released under CC BY-NC-ND 4.0 to protect
% the integrity of the theoretical work during ongoing academic development.
%
% See LICENSE.md for full license text.
%
% File: research_tracks/T3_ALPHA/vcw_exact_minimum.tex
%
% Purpose: Document the minimum of the exact Coleman-Weinberg potential
%   V_CW(n) = n^2 - N_eff * ln(2*sinh(pi*n))
%   and investigate whether combined potentials can produce n*=137.
%
% Numerical results generated by: tools/vcw_minimum_scan.py (2026-05-10)
%
% Background: cw_determinant_full_derivation.tex, nlogn_mechanism.tex

% UBT-AI-PROVENANCE-BEGIN
% schema: ubt-ai-provenance/v1
% tier: C_working
% ai_assistance: disclosed
% human_review: risk-based
% editorial_responsibility: Ing. David Jaroš
% policy: ../../AI_PROVENANCE.md
% notice: Working material; exhaustive human review is not claimed.
% UBT-AI-PROVENANCE-END

\ifdefined\INCLUDEMODE
\else
  \documentclass[12pt,a4paper]{article}
  \usepackage[a4paper, margin=2.5cm]{geometry}
  \usepackage{amsmath, amssymb, amsthm}
  \usepackage{mathtools}
  \usepackage{hyperref}
  \usepackage{xcolor}
  \usepackage{booktabs}
  \usepackage{longtable}
  \usepackage{mdframed}
  \usepackage{tcolorbox}
  \tcbuselibrary{skins,breakable}

  \newtheorem{theorem}{Theorem}[section]
  \newtheorem{lemma}[theorem]{Lemma}
  \newtheorem{proposition}[theorem]{Proposition}
  \theoremstyle{definition}
  \newtheorem{definition}[theorem]{Definition}
  \theoremstyle{remark}
  \newtheorem{remark}[theorem]{Remark}

  \newmdenv[backgroundcolor=yellow!15, linecolor=orange!70, linewidth=1.5pt]{gapbox}
  \newmdenv[backgroundcolor=green!12, linecolor=green!60!black, linewidth=1.5pt]{resultbox}
  \newmdenv[backgroundcolor=blue!6, linecolor=blue!45, linewidth=1.5pt]{insightbox}
  \newmdenv[backgroundcolor=red!8, linecolor=red!60, linewidth=1.5pt]{warnbox}

  \newtcolorbox{statusbox}[1][]{colback=gray!8!white,colframe=black!60,
    arc=3pt,boxrule=1pt,fonttitle=\bfseries,title=Status Summary,#1}

  \newcommand{\Veff}{V_{\mathrm{eff}}}
  \newcommand{\Neff}{N_{\mathrm{eff}}}
  \newcommand{\VCW}{V_{\mathrm{CW}}}

  \title{\textbf{Exact $V_{\mathrm{CW}}$ Minimum and Path to $n^*=137$}\\[4pt]
         \large Track T3\_ALPHA: Approach C --- Coleman-Weinberg analysis}
  \author{Ing.\ David Jaro\v{s}}
  \date{2026-05-10}

  \begin{document}
  \maketitle
\fi

%──────────────────────────────────────────────────────────────────────────────
\section{The Exact Coleman-Weinberg Potential}
\label{sec:vcw:exact}
%──────────────────────────────────────────────────────────────────────────────

The one-loop Coleman-Weinberg functional determinant for $\Neff$ bosonic modes
on $S^1_\psi$ with radius $R_\psi$ (from
\texttt{cw\_determinant\_full\_derivation.tex}) gives the exact effective potential:
\begin{equation}
  \VCW(n) = n^2 - \Neff\ln\!\left(2\sinh(\pi n)\right) .
  \label{eq:VCW}
\end{equation}
The stationarity condition is:
\begin{equation}
  \frac{d\VCW}{dn} = 2n - \Neff\,\pi\coth(\pi n) = 0 .
  \label{eq:stationarity}
\end{equation}
For large $n$ (which covers $n=137$), $\coth(\pi n) \approx 1$, giving the
approximate minimum:
\begin{equation}
  n^*(\VCW) \approx \frac{\Neff\,\pi}{2} .
  \label{eq:nstar_approx}
\end{equation}

\subsection{Key observation: $n^*$ scales linearly with $\Neff$}

The minimum of $\VCW$ is linear in $\Neff$:
\begin{align}
  \Neff = 3 &\Rightarrow n^* \approx 4.71 , \\
  \Neff = 12 &\Rightarrow n^* \approx 18.85 \quad (\text{minimum at } n=19, \text{ prime}) , \\
  \Neff = 87 &\Rightarrow n^* \approx 136.7 \quad (\text{minimum at } n=137, \text{ prime}) .
\end{align}

%──────────────────────────────────────────────────────────────────────────────
\section{Numerical Results}
\label{sec:vcw:numerical}
%──────────────────────────────────────────────────────────────────────────────

\subsection{Minimum of $\VCW$ for various $\Neff$}

Results from \texttt{tools/vcw\_minimum\_scan.py}:

\begin{center}
\begin{tabular}{rrrrr}
  \toprule
  $\Neff$ & $n^*(\VCW)$ & $V_{\mathrm{CW}}(n^*)$ & $n^*$ prime? & Approx.\ $\Neff\pi/2$ \\
  \midrule
  3   &   4.71  & $-22.21$    & YES & 4.71 \\
  6   &   9.42  & $-88.83$    & NO  & 9.42 \\
  12  &  18.85  & $-355.31$   & YES (19) & 18.85 \\
  24  &  37.70  & $-1421.22$  & NO  & 37.70 \\
  48  &  75.40  & $-5684.89$  & NO  & 75.40 \\
  87  & 136.66  & $-18675.76$ & YES (137) & 136.66 \\
  88  & 138.23  & $-19107.55$ & NO  & 138.23 \\
  \bottomrule
\end{tabular}
\end{center}

\subsection{$\Neff$ required for $n^* = 137$}

From $2n^* = \Neff\,\pi\coth(\pi n^*)$ at $n^* = 137$:
\begin{equation}
  \Neff^{\mathrm{needed}} = \frac{2 \times 137}{\pi \coth(\pi \times 137)}
  = \frac{274}{\pi} \approx 87.217 .
  \label{eq:Neff_needed}
\end{equation}
The UBT theory gives $\Neff \in \{3, 12\}$ (from the biquaternion algebra).
The gap is:
\begin{equation}
  \frac{\Neff^{\mathrm{needed}}}{\Neff^{\mathrm{UBT}}}
  = \frac{87.2}{12} \approx 7.27 .
  \label{eq:gap_factor}
\end{equation}

\begin{warnbox}
\textbf{Result}: $\VCW$ alone, with $\Neff = 12$ (the UBT value from
$\mathbb{C}\otimes\mathbb{H}$ algebra), places the minimum at $n^* \approx 19$,
not $n^* = 137$.  Reaching $n^*=137$ from $\VCW$ alone would require
$\Neff \approx 87$, which is not available in UBT.  The gap factor is
$\approx 7.27$.
\end{warnbox}

\subsection{$\VCW$ scan, $n \in [1,200]$, $\Neff=12$ (selected entries)}

\begin{center}
\begin{tabular}{rrcc}
  \toprule
  $n$ & $\VCW(n)$ & Prime? & Local min? \\
  \midrule
  17 & $-351.88$ & YES & NO \\
  18 & $-354.58$ & NO  & NO \\
  \textbf{19} & $\mathbf{-355.28}$ & \textbf{YES} & \textbf{YES (global)} \\
  20 & $-353.98$ & NO  & NO \\
  21 & $-350.68$ & NO  & NO \\
  $\vdots$ & $\vdots$ & & \\
  137 & $13604.22$ & YES & NO \\
  \bottomrule
\end{tabular}
\end{center}

At $n=137$, $\VCW(137) = 13604.22 \gg \VCW(19) = -355.28$.
The gap to overcome is $\Delta V = 13959.50$.

%──────────────────────────────────────────────────────────────────────────────
\section{Combined Potentials}
\label{sec:vcw:combined}
%──────────────────────────────────────────────────────────────────────────────

\subsection{Modular correction $V_{\mathrm{mod}}(n)$}

For prime $n = p$, the modular volume of $X_0(p)$ suggests a correction term:
\begin{equation}
  V_{\mathrm{mod}}(n) = \VCW(n) - C\,\frac{n+1}{3} \quad \text{(for prime } n \text{)},
  \label{eq:Vmod}
\end{equation}
where $C$ is a coupling derived from $S[\Theta]$.  The correction $-C(n+1)/3$
is larger for larger primes, which can in principle shift the minimum upward.

\subsection{$C$ required to flip minimum from $n=19$ to $n=137$}

Both $n=19$ and $n=137$ are prime, so both receive the modular correction.
The condition $V_{\mathrm{mod}}(137) < V_{\mathrm{mod}}(19)$ requires:
\begin{align}
  \VCW(137) - C\,\frac{138}{3}
  &< \VCW(19) - C\,\frac{20}{3} \notag\\
  C\left(\frac{138-20}{3}\right)
  &> \VCW(137) - \VCW(19) = 13959.50 \notag\\
  C\,\cdot\,\frac{118}{3} &> 13959.50 \notag\\
  C &> 354.9 .
  \label{eq:C_needed}
\end{align}

The natural scale of $C$ from UBT is $C_{\mathrm{nat}} = \Neff\,\pi/3 \approx 12.57$.
The required $C$ exceeds the natural scale by a factor of
\begin{equation}
  \frac{C_{\mathrm{min}}}{C_{\mathrm{nat}}} = \frac{354.9}{12.57} \approx 28.2 .
  \label{eq:C_ratio}
\end{equation}

\begin{gapbox}
\textbf{Finding}: The modular correction $V_{\mathrm{mod}}(n) = \VCW(n) - C(n+1)/3$
is insufficient to place the minimum at $n^*=137$ for any natural value of $C$
derivable from UBT with $\Neff=12$.  The required $C \approx 355$ is $\approx 28$
times the natural scale $C_{\mathrm{nat}} \approx 12.6$.

Even with $C = C_{\mathrm{min}}$, the combined minimum lands at $n^*=79$
(not $137$), due to the modular correction pushing the minimum only partially.
\end{gapbox}

\subsection{What $\Neff$ profile would be needed?}

For $\VCW$ alone to give $n^*=137$, one would need $\Neff \approx 87.2$.
Potential sources of enhancement:
\begin{enumerate}
  \item \textbf{Higher-loop corrections}: Each additional loop may multiply
        $\Neff$ by an order-unity factor.  A factor $87/12 \approx 7.27$ from
        loop corrections would require $\sim 3$ loops each contributing factor
        $\sim 2$, which is not natural.
  \item \textbf{Kac-Moody level}: If the WZW term gives level $k_{\mathrm{KM}} = 7$,
        the effective coupling could be $\Neff \to 7 \times 12 = 84 \approx 87$.
        This is the Gap~G3-k conjecture (see
        \texttt{canonical/alpha/gap\_g3k\_closure\_report.tex}).
  \item \textbf{$T^3$ factor}: The background $T^3 \times S^1_\psi$ may
        contribute additional zero-mode degeneracy.
\end{enumerate}

None of these mechanisms is currently derived from $S[\Theta]$.

%──────────────────────────────────────────────────────────────────────────────
\section{Connection to the $n\ln n$ Potential}
\label{sec:vcw:nlogn}
%──────────────────────────────────────────────────────────────────────────────

The $\Veff(n) = n^2 - Bn\ln n$ form (the ``$n\ln n$ potential'', Track~A\_PRIME)
differs fundamentally from $\VCW(n) = n^2 - \Neff\ln(2\sinh(\pi n))$:

\begin{center}
\begin{tabular}{lll}
  \toprule
  Property & $\VCW$ & $\Veff$ ($n\ln n$) \\
  \midrule
  Large-$n$ behaviour & $n^2 - \Neff\,\pi n$ (linear) & $n^2 - Bn\ln n$ (quasi-log) \\
  Minimum location & $n^* \approx \Neff\pi/2$ & $n^* \approx \mathrm{Lambert}(B)$ \\
  $n^*=137$ requires & $\Neff \approx 87$ & $B \approx 46.3$ \\
  UBT value & $\Neff = 12$ & $B_0 = 8\pi \approx 25.1$ \\
  Gap & factor $7.27$ & factor $1.84$ \\
  \bottomrule
\end{tabular}
\end{center}

The $n\ln n$ potential has a smaller gap ($B_0/B_{\mathrm{phenom}} \approx 1.84$
vs.\ $\Neff/\Neff_{\mathrm{needed}} \approx 7.27$), making it the better
starting point for the alpha derivation.  However, the exact $\VCW$ form
does not generate the $n\ln n$ shape (see \texttt{nlogn\_mechanism.tex}).

%──────────────────────────────────────────────────────────────────────────────
\section{Summary}
\label{sec:vcw:summary}
%──────────────────────────────────────────────────────────────────────────────

\begin{center}
\begin{tabular}{lll}
  \toprule
  Result & Value & Classification \\
  \midrule
  $\VCW$ min (exact, $\Neff=12$) & $n^* = 19$ (prime) & [L1] \\
  $\VCW$ min (exact, $\Neff=3$)  & $n^* \approx 4.7$  & [L1] \\
  $\Neff$ for $n^*=137$ ($\VCW$) & $87.2$ & [L0] calculation \\
  Gap factor $87.2/12$           & $7.27$ & [L0] \\
  $C$ needed for $V_{\mathrm{mod}}$, $n^*=137$ & $354.9$ & [L0] calc \\
  $C_{\mathrm{nat}}$ from UBT    & $12.6$ & [MC] \\
  Ratio $C/C_{\mathrm{nat}}$     & $28.2$ & [L0] \\
  $V_{\mathrm{combined}}$ min (using $C=354.9$) & $n^*=79$ & [L0] \\
  \bottomrule
\end{tabular}
\end{center}

\begin{statusbox}
\textbf{$V_{\mathrm{CW}}$ minimum ($\Neff=12$)}: $n^* = 19$ (prime).

\textbf{$V_{\mathrm{combined}}$ minimum}: $n^* = 79$ with $C = 354.9$ (not natural).

\textbf{Potřebný přídavný mechanismus}:
\begin{itemize}
  \item $\Neff^{\mathrm{eff}} \approx 87$ from higher loops or Kac-Moody level $k\approx 7$.
  \item OR: replace $\VCW$ with $\Veff = n^2 - Bn\ln n$ (smaller gap, factor 1.84).
  \item OR: identify physical source of $n\ln n$ term beyond the direct CW determinant.
\end{itemize}

\textbf{Klasifikace}: [L1] pro výsledky o minimech $\VCW$;
[OPEN] pro mechanismus posunu minima na $n^*=137$ bez volného parametru.

\textbf{Závěr}: $\VCW$ sama o sobě nedává $n^*=137$ pro žádnou přirozenou hodnotu
$\Neff$ nebo $C$ odvoditelnou z UBT.  Žádná rozumná kombinace $\VCW +
V_{\mathrm{divisor}} + V_{\mathrm{mod}}$ s přirozenými konstantami nedosáhne
$n^*=137$.  Přístup C \emph{sám o sobě} nemůže uzavřít gap G137-B.
\end{statusbox}

\ifdefined\INCLUDEMODE
\else
  \end{document}
\fi
